% ============================================================
% NEW §IV-E (revised): Deployability Safety Standards
% Replaces current §IV-E with risk-based reasoning as primary basis
% ============================================================

\subsection{Deployability Safety Standard}
\label{sec:deployability_standard}

\subsubsection{Risk-based derivation of DR@5s threshold}

We derive the DR@$5\,\text{s}$ threshold from first principles of UAV
power-line inspection risk.
A typical inspection UAV cruises at $v \approx 5\,\text{m/s}$.
Under the Chinese national standard \emph{GB\,26860-2011},
the minimum safe approach distance to a $\SI{500}{\kV}$ transmission line is
$d_{\min} = \SI{5}{\metre}$.
Over a $5\,\text{s}$ response window, a spoofed UAV drifts up to
\[
  \Delta x = v \cdot \Delta t = 5 \times 5 = \SI{25}{\metre},
\]
leaving $25 - 5 = \SI{20}{\metre}$ of margin for the operator to abort
the approach after a timely alarm.
If DR@5s~$\geq 0.85$, at least 85\% of attacks are detected before the UAV
exhausts this safety margin; we adopt this as the primary deployability threshold.

The analogous threshold adopted in human-machine interface engineering
\cite{eemua191} and process-control alarm management \cite{iso11064}
corroborates the $5\,\text{s}$ response budget, but we treat these domain
analogies as supporting evidence rather than primary justification.

\subsubsection{Risk-based derivation of MTBFA threshold}

A false alarm requires the remote pilot to intervene (verify sensor data,
abort the mission leg, re-arm).
During a typical $\SI{1}{\hour}$ inspection tour the pilot can reasonably
absorb at most 10 nuisance alarms without significant workload increase
\cite{axelsson2000base,iso11064}, giving a minimum tolerable
$\text{MTBFA} \geq 0.1\,\text{h}$.

\subsubsection{Deployability verdict}

A model passes the time-aware deployability standard if and only if both
conditions hold simultaneously:
\begin{equation}
  \text{DR@5s} \geq 0.85 \quad \text{and} \quad \text{MTBFA} \geq 0.1\,\text{h}.
\end{equation}
Table~\ref{tab:deployability} reports the pass/fail verdict for all seven
detectors evaluated on the MATLAB test set.
